\documentclass{beamer}
\usetheme{Madrid}

\usepackage[utf8]{inputenc} % solo si usas pdflatex
\usepackage[T1]{fontenc}    % solo si usas pdflatex

\usepackage{amsmath,amssymb}
\usepackage{microtype}
\usepackage{cancel}         % si lo usas
\usepackage{fontawesome}    % si lo usas
% \usepackage{xypic}        % solo si lo usas
% \usepackage{multicol}     % normalmente no; usar columns de beamer
% \usepackage{amsthm}       % solo si defines teoremas propios
% \usepackage{scrextend}    % solo si lo usas

\usepackage{graphicx}
\graphicspath{{images/}}

\hypersetup{pdfborder={0 0 0}}


\title{Modelos y formatos de almacenamiento de datos}
\author{Rodrigo Blanco Betes. Alfa Formación}
\date{Febrero del 2026}

\begin{document}

\begin{frame}
\titlepage
\end{frame}

\begin{frame}
\tableofcontents
\end{frame}


% --- Slide deck content (Beamer) ---
% Nota: Ignoro encabezados globales (documentclass, theme, etc.) como pediste.
% Incluye esta sección dentro de tu .tex beamer.

\section{Object y block storage}

\begin{frame}{¿Por qué hablamos de tipos de almacenamiento?}
En Big Data y Cloud, el modelo de almacenamiento condiciona:
\begin{itemize}
    \item \textbf{Rendimiento}: latencia y throughput.
    \item \textbf{Escalabilidad}: crecimiento y distribución.
    \item \textbf{Coste}: capacidad, tráfico, operación.
    \item \textbf{Casos de uso}: analítica, transaccional, backups, etc.
\end{itemize}
\end{frame}

\begin{frame}{Block Storage: idea básica}
\textbf{Block Storage} guarda datos en \textbf{bloques} (unidades de tamaño fijo o gestionadas como bloques).
\begin{itemize}
    \item Cada bloque tiene un identificador/dirección.
    \item Encima suele haber un \textbf{sistema de ficheros} (ext4, NTFS, XFS, \dots) que organiza archivos y directorios.
    \item Se comporta como un \textbf{disco} (volumen) para el sistema operativo.
\end{itemize}
\end{frame}

\begin{frame}{Block Storage: cuándo conviene}
\textbf{Ventajas}
\begin{itemize}
    \item Baja latencia y buen rendimiento en I/O aleatorio.
    \item Adecuado para workloads \textbf{transaccionales} y bases de datos.
\end{itemize}

\textbf{Limitaciones}
\begin{itemize}
    \item Escalado y compartición pueden ser más complejos (dependen del proveedor/arquitectura).
    \item Normalmente no es el modelo más cómodo para “data lake” masivo.
\end{itemize}
\end{frame}

\begin{frame}{Block Storage en el ecosistema Hadoop}
Aunque Hadoop trabaja a nivel de \textbf{ficheros}, su almacenamiento distribuido usa un enfoque basado en bloques.

\begin{itemize}
    \item \textbf{HDFS} divide los ficheros en \textbf{bloques grandes} (p.ej. 128MB/256MB).
    \item Los bloques se distribuyen entre nodos (\textbf{DataNodes}).
    \item Se replican bloques (típicamente 3) para tolerancia a fallos.
\end{itemize}

\medskip
\textbf{Idea clave:} bloques grandes $\Rightarrow$ paralelismo + throughput alto en lecturas/escrituras secuenciales.
\end{frame}

\begin{frame}{Object Storage: idea básica}
\textbf{Object Storage} guarda datos como \textbf{objetos} en un espacio “plano” (contenedores/buckets).
Cada objeto incluye:
\begin{itemize}
    \item Los \textbf{datos}.
    \item \textbf{Metadatos} (clave-valor) ricos.
    \item Un \textbf{identificador} único (nombre/clave del objeto).
\end{itemize}

Acceso típico: \textbf{API HTTP/REST} (no un sistema de archivos clásico).
\end{frame}

\begin{frame}{Object Storage: cuándo conviene}
\textbf{Ventajas}
\begin{itemize}
    \item Escala masiva y alta durabilidad.
    \item Muy bueno para datos \textbf{no estructurados} y grandes volúmenes.
    \item Metadatos útiles para catalogación y gobierno del dato.
\end{itemize}

\textbf{Limitaciones}
\begin{itemize}
    \item Latencia mayor que block para I/O fino.
    \item Semántica diferente: no siempre hay “rename”/“append” eficientes como en FS tradicionales.
\end{itemize}
\end{frame}

\begin{frame}{Ejemplo: Azure Blob Storage (Object Storage)}
\textbf{Azure Blob Storage} es almacenamiento de objetos en la nube:
\begin{itemize}
    \item Organización en \textbf{contenedores} y \textbf{blobs} (objetos).
    \item Acceso por \textbf{HTTP/HTTPS} y SDKs.
    \item Integración con analítica: Spark, Synapse, Databricks, etc.
\end{itemize}

\medskip
Uso típico: data lake, backups, logs, multimedia, datasets analíticos.
\end{frame}

\begin{frame}{Hadoop + Cloud: desacoplar cómputo y almacenamiento}
En on-premises clásico:
\begin{itemize}
    \item Cómputo y almacenamiento suelen ir juntos (HDFS en los nodos).
\end{itemize}

En cloud:
\begin{itemize}
    \item Es común usar \textbf{Object Storage} como “data lake”.
    \item Motores (Spark/Hive/Presto) leen directamente de Blob/S3.
    \item Beneficio: \textbf{elasticidad} (escala cómputo sin mover datos).
\end{itemize}
\end{frame}

\begin{frame}{Comparativa rápida}
\begin{itemize}
    \item \textbf{Unidad}: bloques vs objetos.
    \item \textbf{Acceso}: como disco/sistema de archivos vs API (REST).
    \item \textbf{Rendimiento típico}: block (baja latencia) vs object (alto throughput masivo).
    \item \textbf{Casos de uso}: BD/VMs vs data lake/backups/datos no estructurados.
    \item \textbf{Ejemplos}: HDFS (bloques), Azure Managed Disks (block) vs Azure Blob (object).
\end{itemize}
\end{frame}

\begin{frame}{Mensaje final}
\begin{itemize}
    \item \textbf{Block Storage}: “disco” para I/O fino y transaccional.
    \item \textbf{Object Storage}: “repositorio masivo” para datos escalables y analítica.
    \item En Big Data moderno: tendencia a \textbf{object storage + motores de cómputo desacoplados}.
\end{itemize}
\end{frame}

%-----------------------------------

%------------------------------------------------
\begin{frame}{¿Object Storage es realmente diferente o solo marketing?}
\textbf{Pregunta capciosa:}

Si todo termina siendo bloques en discos…  
¿no es lo mismo que Block Storage?

\medskip

\textbf{Idea:}
\begin{itemize}
    \item A nivel \textbf{físico}: no es radicalmente distinto.
    \item A nivel \textbf{arquitectónico y de abstracción}: sí es muy distinto.
\end{itemize}
\end{frame}
%------------------------------------------------

\begin{frame}{Nivel físico: lo que no cambia}
En ambos casos:
\begin{itemize}
    \item Hay discos (HDD/SSD).
    \item Hay bloques físicos.
    \item Hay servidores en centros de datos.
    \item Hay replicación y redes.
\end{itemize}

\medskip
\textbf{Conclusión:} Todo almacenamiento termina siendo bloques en algún disco.
\end{frame}
%------------------------------------------------

\begin{frame}{Nivel lógico: lo que sí cambia}
\textbf{Block Storage}
\begin{itemize}
    \item Se comporta como un disco.
    \item Tú gestionas sistema de archivos.
    \item Permite modificar partes pequeñas.
\end{itemize}

\medskip

\textbf{Object Storage}
\begin{itemize}
    \item Se comporta como un servicio.
    \item No hay sistema de archivos real.
    \item Se gestionan objetos completos.
    \item Acceso vía API (HTTP).
\end{itemize}
\end{frame}
%------------------------------------------------

\begin{frame}{La diferencia real está en la arquitectura}
\begin{itemize}
    \item Object Storage está diseñado para:
    \begin{itemize}
        \item Escalar masivamente.
        \item Distribuir datos automáticamente.
        \item Replicar sin intervención del usuario.
        \item Integrarse como servicio en arquitecturas cloud.
    \end{itemize}
    \item Block Storage está diseñado para:
    \begin{itemize}
        \item Simular un disco.
        \item Soportar bases de datos y máquinas virtuales.
    \end{itemize}
\end{itemize}
\end{frame}

% --- Continuación (Beamer) ---
% Nota: ignoro encabezados globales y mantengo el estilo de frames.

\section{Formatos: Avro y Parquet}

\begin{frame}{Formatos de ficheros en Hadoop (visión rápida)}
En el ecosistema Hadoop/Spark aparecen varios formatos comunes:
\begin{itemize}
    \item \textbf{Text}: CSV/TSV/JSON/logs (simple, pero poco eficiente).
    \item \textbf{SequenceFile}: binario clave-valor (históricamente ligado a MapReduce).
    \item \textbf{Avro}: binario orientado a filas, \textbf{basado en esquema}.
    \item \textbf{Columnar}: Parquet/ORC (optimizado para analítica).
\end{itemize}
\end{frame}

\begin{frame}{Apache Avro: ¿qué es?}
\textbf{Avro} es un sistema de \textbf{serialización} con \textbf{esquema} (schema) usado en Hadoop.

\begin{itemize}
    \item Los datos se almacenan en binario compacto.
    \item El \textbf{esquema} (en JSON) describe la estructura.
    \item Permite \textbf{evolución de esquemas} (añadir campos, compatibilidad hacia atrás, etc.).
\end{itemize}

\medskip
Uso típico: intercambio de datos entre sistemas, ingestión, streaming y etapas intermedias.
\end{frame}

\begin{frame}{Avro: ventajas clave}
\begin{itemize}
    \item \textbf{Compacto} y eficiente en red/almacenamiento.
    \item \textbf{Autodescriptivo}: el esquema puede viajar con los datos (dependiendo del contenedor).
    \item \textbf{Evolución de esquema}: muy útil cuando los datos cambian con el tiempo.
    \item Buen encaje con pipelines (p.ej. Kafka $\rightarrow$ procesamiento $\rightarrow$ data lake).
\end{itemize}
\end{frame}

\begin{frame}{Avro: tipos de datos (básicos y complejos)}
\textbf{Básicos:}
\begin{itemize}
    \item \texttt{null}, \texttt{boolean}, \texttt{int}, \texttt{long}, \texttt{float}, \texttt{double}
    \item \texttt{bytes}, \texttt{string}
\end{itemize}

\textbf{Complejos:}
\begin{itemize}
    \item \texttt{record} (estructura con campos)
    \item \texttt{enum}, \texttt{array}, \texttt{map}
    \item \texttt{union} (p.ej. \texttt{["null","string"]} para opcionales)
    \item \texttt{fixed}
\end{itemize}
\end{frame}

\begin{frame}[fragile]{Ejemplo de esquema Avro (JSON)}
\begin{verbatim}
{
  "type": "record",
  "name": "Usuario",
  "namespace": "ejemplo",
  "fields": [
    {"name": "id", "type": "int"},
    {"name": "nombre", "type": "string"},
    {"name": "email", "type": ["null","string"], "default": null}
  ]
}
\end{verbatim}

\medskip
\textbf{Idea}: \texttt{union} permite valores nulos (campo opcional).
\end{frame}

\begin{frame}{Relación con Java: Avro en la práctica}
Avro nació con fuerte integración en el ecosistema \textbf{Java}.

\begin{itemize}
    \item Flujo típico en Java:
    \begin{itemize}
        \item Definir el esquema (\texttt{.avsc}) o un IDL.
        \item \textbf{Generar clases Java} (POJOs) a partir del esquema.
        \item Serializar/deserializar con las librerías de Avro.
    \end{itemize}
    \item Ventaja: \textbf{tipado} y modelos de datos consistentes entre microservicios y batch.
\end{itemize}

\medskip
En entornos Hadoop, esto facilita escribir jobs (MapReduce) o componentes de ingestión en Java.
\end{frame}

\begin{frame}{Parquet: formato columnar para analítica}
\textbf{Apache Parquet} es un formato \textbf{columnar} pensado para consultas analíticas.

\begin{itemize}
    \item Almacena datos \textbf{por columnas} (no por filas).
    \item Muy eficiente cuando se consultan \textbf{solo algunas columnas}.
    \item Permite \textbf{compresión} y \textbf{codificación} por columna.
\end{itemize}

\medskip
Uso típico: data lakes, tablas analíticas (Spark SQL, Hive, Impala, Trino/Presto).
\end{frame}

\begin{frame}{¿Qué aporta lo columnar?}
Cuando preguntas: \emph{``dame media(importe) por país''}:

\begin{itemize}
    \item En formato fila (CSV/Avro): lees muchas columnas aunque no las uses.
    \item En \textbf{Parquet}: puedes leer solo \textbf{importe} y \textbf{país}.
\end{itemize}

\medskip
Esto mejora:
\begin{itemize}
    \item \textbf{Tiempo de lectura} (menos I/O).
    \item \textbf{Coste} en cloud (menos datos escaneados).
    \item \textbf{Compresión} (valores similares por columna).
\end{itemize}
\end{frame}

\begin{frame}{Relación con Hadoop/Spark}
\begin{itemize}
    \item Parquet es un estándar de facto en el \textbf{ecosistema Hadoop} moderno.
    \item Motores como Spark y Hive lo optimizan con:
    \begin{itemize}
        \item \textbf{Predicate pushdown} (filtrar lo antes posible).
        \item \textbf{Column pruning} (leer solo columnas necesarias).
        \item Estadísticas y metadatos para planificar mejor consultas.
    \end{itemize}
\end{itemize}
\end{frame}

\begin{frame}{Relación con Java: Parquet en la práctica}
También hay integración directa con \textbf{Java} (librerías Apache Parquet).

\begin{itemize}
    \item Casos típicos en Java:
    \begin{itemize}
        \item Generar datasets Parquet en pipelines batch.
        \item Escribir/leer Parquet usando esquemas (a menudo con Avro como modelo).
    \end{itemize}
    \item Patrón común: \textbf{Avro schema} $\rightarrow$ escribir Parquet (muy usado en pipelines).
\end{itemize}

\medskip
\textbf{Idea}: Avro facilita el modelo y compatibilidad; Parquet optimiza almacenamiento y consulta.
\end{frame}

\begin{frame}{Avro vs Parquet: ¿cuándo usar cada uno?}
\begin{itemize}
    \item \textbf{Avro} (fila + esquema):
    \begin{itemize}
        \item Ingestión, intercambio, streaming.
        \item Escritura secuencial y evolución de esquema frecuente.
    \end{itemize}
    \item \textbf{Parquet} (columna):
    \begin{itemize}
        \item Data lake analítico, BI, agregaciones, consultas repetidas.
        \item Rendimiento y coste optimizados en lectura selectiva.
    \end{itemize}
\end{itemize}
\end{frame}

\begin{frame}{Pipeline típico en Big Data (con Java en la cadena)}
\begin{itemize}
    \item \textbf{Java} (producer/servicio) genera eventos con esquema Avro.
    \item Se almacenan o transportan (p.ej. en colas/streams) manteniendo el esquema.
    \item Proceso batch/streaming transforma y persiste en \textbf{Parquet} para analítica.
\end{itemize}

\medskip
Resultado: \textbf{compatibilidad} y \textbf{gobierno del dato} (Avro) + \textbf{rendimiento analítico} (Parquet).
\end{frame}


%------------------------------------------------
\section{Colas, ZooKeeper y Apache Kafka}
%------------------------------------------------

\begin{frame}{Introducción a las colas}
Las \textbf{colas de mensajes} permiten desacoplar productores y consumidores.

\begin{itemize}
    \item Un \textbf{productor} envía mensajes a una cola.
    \item Uno o varios \textbf{consumidores} los procesan.
    \item Comunicación asíncrona.
\end{itemize}

\medskip
Objetivo principal: \textbf{desacoplar sistemas} y mejorar escalabilidad y tolerancia a fallos.
\end{frame}

\begin{frame}{¿Por qué usar colas?}
\begin{itemize}
    \item \textbf{Desacoplamiento}: productor y consumidor no necesitan conocerse.
    \item \textbf{Escalabilidad}: múltiples consumidores pueden procesar en paralelo.
    \item \textbf{Tolerancia a fallos}: los mensajes se almacenan hasta procesarse.
    \item \textbf{Asincronía}: mejora rendimiento en sistemas distribuidos.
\end{itemize}

\medskip
Muy utilizadas en microservicios, Big Data y sistemas event-driven.
\end{frame}

%------------------------------------------------
\subsection*{ZooKeeper}
%------------------------------------------------

\begin{frame}{ZooKeeper: ¿qué es?}
\textbf{Apache ZooKeeper} es un servicio distribuido de \textbf{coordinación}.

No es una cola ni una base de datos tradicional.

Proporciona:
\begin{itemize}
    \item Configuración centralizada.
    \item Sincronización distribuida.
    \item Gestión de líderes.
    \item Registro de servicios.
\end{itemize}
\end{frame}

\begin{frame}{ZooKeeper: arquitectura}
\begin{itemize}
    \item Funciona como un \textbf{ensemble} (conjunto) de nodos.
    \item Uno actúa como \textbf{líder}, los demás como seguidores.
    \item Utiliza un modelo jerárquico tipo sistema de archivos (\texttt{znodes}).
\end{itemize}

\medskip
Principio clave: \textbf{consistencia fuerte} (protocolo tipo quorum).
\end{frame}

\begin{frame}{ZooKeeper: propiedades}
\begin{itemize}
    \item \textbf{Alta disponibilidad} mediante replicación.
    \item \textbf{Consistencia} garantizada.
    \item Datos pequeños y en memoria (no pensado para grandes volúmenes).
    \item Operaciones rápidas (lecturas frecuentes).
\end{itemize}

\medskip
En el ecosistema Hadoop, se usa para coordinar servicios como HBase o Kafka (versiones clásicas).
\end{frame}

%------------------------------------------------
\subsection*{Apache Kafka}
%------------------------------------------------

\begin{frame}{Apache Kafka: visión general}
\textbf{Apache Kafka} es un \textbf{broker de mensajería distribuido} orientado a eventos.

Modelo:
\begin{itemize}
    \item Productores publican mensajes.
    \item Mensajes se almacenan en \textbf{topics}.
    \item Consumidores leen desde esos topics.
\end{itemize}

Diseñado para alto throughput y procesamiento en tiempo real.
\end{frame}

\begin{frame}{Kafka: arquitectura básica}
\begin{itemize}
    \item \textbf{Brokers}: servidores que almacenan datos.
    \item \textbf{Topics}: categorías de mensajes.
    \item \textbf{Partitions}: subdivisiones de un topic para paralelismo.
    \item \textbf{Consumers Groups}: permiten escalado horizontal.
\end{itemize}

\medskip
Los datos se almacenan en disco de forma secuencial (muy eficiente).
\end{frame}

\begin{frame}{Kafka: características principales}
\begin{itemize}
    \item \textbf{Alto rendimiento} (millones de mensajes/segundo).
    \item \textbf{Persistencia en disco}.
    \item \textbf{Escalabilidad horizontal}.
    \item \textbf{Replicación} entre brokers.
    \item Garantías configurables: \textit{at-most-once}, \textit{at-least-once}, \textit{exactly-once}.
\end{itemize}
\end{frame}

\begin{frame}{Kafka y Java}
Kafka está desarrollado en \textbf{Java/Scala} y ofrece API nativa en Java.

Uso típico en Java:
\begin{itemize}
    \item Crear un \textbf{Producer} para enviar eventos.
    \item Crear un \textbf{Consumer} para procesarlos.
    \item Integración con microservicios Spring Boot.
\end{itemize}

\medskip
Muy común en pipelines:
\[
Java\ Service \rightarrow Kafka \rightarrow Spark/Flink \rightarrow Data\ Lake
\]
\end{frame}

\begin{frame}{Relación ZooKeeper -- Kafka}
\begin{itemize}
    \item Versiones clásicas de Kafka usan \textbf{ZooKeeper} para:
    \begin{itemize}
        \item Elección de líder.
        \item Gestión de metadatos del clúster.
        \item Coordinación entre brokers.
    \end{itemize}
    \item Versiones modernas pueden funcionar sin ZooKeeper (modo KRaft).
\end{itemize}

\medskip
Idea clave: ZooKeeper coordina; Kafka transporta eventos.
\end{frame}

\begin{frame}{Conclusión sobre Zookeper y Kafka}
\begin{itemize}
    \item \textbf{Colas}: desacoplamiento y comunicación asíncrona.
    \item \textbf{ZooKeeper}: coordinación distribuida y consistencia.
    \item \textbf{Kafka}: broker de mensajería distribuido y escalable.
\end{itemize}

Base fundamental de arquitecturas modernas \textbf{event-driven} y Big Data.
\end{frame}

%------------------------------------------------
\section{Bases de datos relacionales vs NoSQL}
%------------------------------------------------

\begin{frame}{Modelo Relacional vs NoSQL}
Dos grandes enfoques en almacenamiento de datos:

\begin{itemize}
    \item \textbf{Bases de datos relacionales (SQL)}
    \item \textbf{Bases de datos NoSQL}
\end{itemize}

La elección depende del tipo de datos, escalabilidad y requisitos del sistema.
\end{frame}

%------------------------------------------------
\begin{frame}{Bases de datos relacionales}
Modelo basado en:
\begin{itemize}
    \item \textbf{Tablas} (filas y columnas).
    \item \textbf{Esquema fijo}.
    \item Relaciones mediante \textbf{claves primarias y foráneas}.
    \item Lenguaje estándar: \textbf{SQL}.
\end{itemize}

\medskip
Ejemplos: MySQL, PostgreSQL, Oracle, SQL Server.
\end{frame}

\begin{frame}{Propiedades del modelo relacional}
\begin{itemize}
    \item Cumplen propiedades \textbf{ACID}.
    \item Alta consistencia.
    \item Adecuadas para sistemas transaccionales (OLTP).
    \item Escalado tradicionalmente vertical.
\end{itemize}

\medskip
Muy utilizadas en banca, ERP, e-commerce tradicional.
\end{frame}

%------------------------------------------------
\begin{frame}{¿Qué es NoSQL?}
\textbf{NoSQL} = “Not Only SQL”.

Diseñadas para:
\begin{itemize}
    \item Escalabilidad horizontal.
    \item Datos semiestructurados o no estructurados.
    \item Grandes volúmenes (Big Data).
\end{itemize}

\medskip
No siguen necesariamente el modelo tabular clásico.
\end{frame}

%------------------------------------------------
\begin{frame}{Tipos de bases de datos NoSQL}
Cuatro categorías principales:

\begin{itemize}
    \item \textbf{Clave-valor}
    \item \textbf{Orientadas a documentos}
    \item \textbf{Orientadas a columnas}
    \item \textbf{De grafos}
\end{itemize}
\end{frame}

%------------------------------------------------
\begin{frame}{1. Clave-Valor}
Modelo más simple:

\[
Clave \rightarrow Valor
\]

\begin{itemize}
    \item Muy rápidas.
    \item Escalables.
    \item El valor es opaco para el sistema.
\end{itemize}

\medskip
Ejemplos: Redis, Riak.
\end{frame}

%------------------------------------------------
\begin{frame}{2. Orientadas a Documentos}
Almacenan documentos (normalmente JSON/BSON).

\begin{itemize}
    \item Estructura flexible (sin esquema rígido).
    \item Documentos auto-contenidos.
    \item Consultas sobre campos internos.
\end{itemize}

\medskip
Ejemplos: MongoDB, CouchDB.
\end{frame}

%------------------------------------------------
\begin{frame}{3. Orientadas a Columnas}
Inspiradas en BigTable (Google).

\begin{itemize}
    \item Almacenan datos por familias de columnas.
    \item Muy eficientes en grandes volúmenes.
    \item Buen rendimiento en escrituras masivas.
\end{itemize}

\medskip
Ejemplos: Cassandra, HBase.
\end{frame}

%------------------------------------------------
\begin{frame}{4. Bases de datos de Grafos}
Modelo basado en:

\begin{itemize}
    \item \textbf{Nodos}
    \item \textbf{Relaciones}
    \item \textbf{Propiedades}
\end{itemize}

Optimizado para consultas sobre relaciones complejas.

\medskip
Ejemplos: Neo4j.
\end{frame}

%------------------------------------------------
\begin{frame}{ACID vs BASE}
\textbf{Relacionales:}
\begin{itemize}
    \item ACID (Atomicidad, Consistencia, Aislamiento, Durabilidad)
\end{itemize}

\textbf{Muchas NoSQL:}
\begin{itemize}
    \item BASE (Basically Available, Soft state, Eventually consistent)
\end{itemize}

\medskip
Trade-off entre consistencia fuerte y escalabilidad.
\end{frame}

%------------------------------------------------
\begin{frame}{Comparación general}
\begin{itemize}
    \item \textbf{Relacional}: estructura rígida, integridad fuerte, SQL estándar.
    \item \textbf{NoSQL}: flexibilidad, escalabilidad horizontal, modelos variados.
    \item Relacional $\rightarrow$ transacciones críticas.
    \item NoSQL $\rightarrow$ Big Data, sistemas distribuidos, microservicios.
\end{itemize}
\end{frame}

%------------------------------------------------
\begin{frame}{Conclusión final}
No existe una solución universal.

\begin{itemize}
    \item Elegir según:
    \begin{itemize}
        \item Tipo de datos.
        \item Volumen.
        \item Necesidades de consistencia.
        \item Escalabilidad.
    \end{itemize}
\end{itemize}

En arquitecturas modernas es común combinar ambos enfoques.
\end{frame}


\end{document}